% ==============================================================================
% РАЗДЕЛ 2.5 (DISS-005-B1350 .. DISS-005-B1392)
% Практическое применение в промышленности и выводы по КР
% ==============================================================================

% DISS-005-B1350
\subsection{Практическое применение в промышленности}
\label{sec:ch02_kr_industrial_application}

% DISS-005-B1351
С 2023 года было введено в эксплуатацию 6 терминалов БАВР и 2 терминала РЗА с подключением токовых каналов с помощью катушек Роговского\commentnote{DISS-005-C0088}{Алексей Евдаков [2]}{Можно указать количество. Вроде бы для 0,4кВ такой один, а для 6-35 -- больше. Могу глянуть и сказать более точно для каждого примера}\commentnote{DISS-005-C0089}{Евдаков Алексей}{С ИПМ (КР+ёмксотн.дел.)\\ 209 -- БАВР 0,4кВ\\ 331 -- БАВР 6кВ (Новосиб)\\ 345 -- БАВР 6кВ + УКРМ (Новосиб)\\ 363 -- БАВР 6кВ (Павлик)\\ + ещё (отгружены но не налажены 556, 559, 560, собраны и будут отгружены 717, 751, 769, 770, 771)\\ \\ КР + ТН\\ 255 -- БАВР 6кВ (Тюмень)\\ 361 -- РЗА 6кВ (Тюмень)\\ 362 -- РЗА 6кВ (Тюмень)\\ 458 -- БАВР (НС Крымская МВК)\\ \\ + ещё (отгружены но не налажены 373, 549, 550, 551, 552, 572 собраны и будут отгружены 736, 737)}\commentnote{DISS-005-C0090}{Евдаков Алексей}{То есть\\ Введено\\ 6 БАВР\\ 2 РЗА\\ \\ И ещё 16 БАВР ожидают своей наладки.}. Данные решения активно применялись в следующих вариантах:

% DISS-005-B1352
При модернизации существующих подстанций 6~$\div$~35~кВ с целью удешевления и упрощения осуществления монтажных работ за счёт отсутствия необходимости в подключении токовых цепей.

% DISS-005-B1353
При установки современных терминалов БАВР, предназначенных для работы на подстанциях 0,4~кВ. Это обусловлено полным отсутствием каких-либо измерительных датчиков тока.

% DISS-005-B1354
При установке БАВР на современных подстанциях, выполненных на ячейках, осуществляющих полный перечень защит, и, следовательно, исключающих возможность подключения дополнительных цепей. Отметим, что в этом случае, дополнительно был разработан промежуточный преобразователь аналоговых сигналов, осуществляющий оцифровку и передачу токов и напряжений через SV (МЭК-61850).

% DISS-005-B1355
Во всех случаях был осуществлён комплекс пуско-наладочных работ, и с того времени устройства показали успешное функционирование не хуже классических устройств на базе ТТ. Опытно-промышленные испытания завершены в конце 2023 года, подход рекомендован к внедрению.

% DISS-005-B1356
Далее приведён пример осциллограммы срабатывания БАВР в сети 6~кВ с применением КР (Рисунок~\ref{fig:ch02_bavr_kr_tt_oscillograms}а). Здесь ОЗЗ перешло в междуфазное КЗ (тока t1), и ввиду просадки линейных напряжений отработал БАВР (точка t2). Время реакции пускового органа составило 6~мс (t1–>t2), а время возврата РНМ 4~мс. На второй секции наблюдается небольшой бросок тока, однако без апериодической составляющей. Для примера приведено аналогичное срабатывание при развитии ОЗЗ для другого объекта (Рисунок~\ref{fig:ch02_bavr_kr_tt_oscillograms}б) с классическими ТТ. Как можем видеть, при включении секционного выключателя наблюдается бросок тока с апериодической составляющей ввиду чего, происходит насыщение ТТ.

\begin{figure}[htbp]
  \centering
  \begin{tabular}{cc}
    % DISS-005-B1357
    \includegraphics[width=0.48\textwidth]{image77.png} &
    % DISS-005-B1358
    \includegraphics[width=0.48\textwidth]{image78.png} \\
    % DISS-005-B1359
    а)\commentnote{DISS-005-C0091}{Евдаков Алексей}{t00331\_231204\_1942\_\_0040 «Эксплуатация\textbackslash 19.01.2024\textbackslash С компа\textbackslash osc»} &
    % DISS-005-B1360
    б)\commentnote{DISS-005-C0092}{Евдаков Алексей}{t00210\_240301\_1259\_\_0021 Буфер (Алексей)\textbackslash Отчёты\textbackslash ТЭЛ Поволжье. г. Нижний Новгород\textbackslash osc}
  \end{tabular}
  % DISS-005-B1361
  % DISS-005-B1362
  \caption{Срабатывание БАВР, t1 -- начало аварии, t2 -- момент срабатывания БАВР\\
  а) с применением КР, б) с применением ТТ*.\\
  {\footnotesize *Примечание: классическое срабатывание приведено с другого объекта}}
  \label{fig:ch02_bavr_kr_tt_oscillograms}
\end{figure}

% DISS-005-B1363
Более того, весной 2024 года были внедрены и опробованы новые реализации БАВР с функциями управления устройств компенсации реактивной мощности (УКРМ), выполненные на батареях статических конденсаторов (БСК). Это позволяет судить о достаточности точности измерения корректного для управления реактивной мощностью.

% DISS-005-B1364
При включении ёмкостной нагрузки наблюдаются сильные броски (с 25А до 200А, Рисунок~\ref{fig:ch02_bavr_kr_ukrm_vnr}а точка t2), которые обусловлены физикой процесса измерения сигнала. Данная особенность измерения сигналов в будущем при большей частоте дискретизации сигналов может быть использована для более быстрого и точного реагирования на особые режимы.

% DISS-005-B1365
При срабатывании функции восстановления нормального режима (ВНР) во всех случаях наблюдались высокочастотные помехи в измерениях тока с КР (Рисунок~\ref{fig:ch02_bavr_kr_ukrm_vnr}б t1–>t2), данный аспект требует дополнительных исследований в будущих работах.

\begin{figure}[htbp]
  \centering
  \begin{tabular}{cc}
    % DISS-005-B1366
    \includegraphics[width=0.48\textwidth]{image79.png} &
    % DISS-005-B1367
    \includegraphics[width=0.48\textwidth]{image80.png} \\
    % DISS-005-B1368
    а)\commentnote{DISS-005-C0093}{Евдаков Алексей}{t00345\_240628\_1658\_\_0008} &
    % DISS-005-B1369
    б)
  \end{tabular}
  % DISS-005-B1370
  % DISS-005-B1371
  % DISS-005-B1372
  \caption{Осциллограммы терминала БАВР с применением КР:\\
  а) при коммутации УКРМ, t1 -- подача команды включения, t2 -- броски ёмкостного тока;\\
  б) при функционировании ВНР, t1 -- включение в параллель, t2 -- отключение параллели.}
  \label{fig:ch02_bavr_kr_ukrm_vnr}
\end{figure}

% DISS-005-B1373
В завершение отметим, что выходным сигналом КР является напряжение, амплитуда которого составляет «мВ», в следствии чего, при прокладке длинного соединяющего провода по кабель-каналам, даже с учётом его экранирования, могут наводиться различные шумы. В сетях 0,4~кВ, указанный эффект не столь существенен, и адекватная точность измерения обеспечивается уже при нескольких амперах -- десятках ампер. Однако в сетях с более высоким напряжением, может оказывать более заметное влияние и требует дальнейшего исследования. Решением этой проблемы может выступать вариант измерений сигнала непосредственно в ячейке промежуточным преобразователем с последующей передачей сигнала по цифровому каналу (SV поток по стандарту МЭК-61850).

% DISS-005-B1374
\subsection{Выводы о применении катушек Роговского для БАВР}
\label{sec:ch02_conclusions_rogowski}

% DISS-005-B1375
Проведенное исследование продемонстрировало актуальность применения катушек Роговского (КР) в устройствах быстродействующего автоматического ввода резерва (БАВР). Основными достоинствами КР являются их более быстрый переход в установившийся режим, минимальное влияние апериодической составляющей, отсутствие насыщения, а также простота монтажа и интеграции в современные устройств.

% DISS-005-B1376
Анализ динамических характеристик показал, что при расчёте первой гармоники (F1) время от начала переходного процесса до установившегося режима (1\% и 1$^\circ$ для класса точности 5Р) для КР не превышает период промышленной частоты (20~мс). В то время как для ТТ, подверженных насыщению и большему влиянию апериодической составляющей, может быть значительно больше. Это обеспечивает более надежное и своевременное срабатывание БАВР, минимизируя последствия аварийных ситуаций.

% DISS-005-B1377
Преимущество КР в точности измерений проявляется при сложных аварийных режимах, таких как короткие замыкания с высокой апериодической составляющей или с большим уровнем тока, приводящем к насыщению ТТ. Исследование работы реле направления мощности (РНМ) подтвердило, что при использовании КР траектории вектора РНМ остаются стабильными, тогда как ТТ могут выдавать ложные срабатывания.

% DISS-005-B1378
Практическое внедрение терминалов БАВР с использованием КР на подстанциях различного уровня напряжения (от 0,4~кВ до 35~кВ) подтвердило эффективность предложенного решения и надёжность функционирования в различных режимах.

% DISS-005-B1379
Следует отметить перспективу одновременного использования сигналов как от классических трансформаторов тока, так и от катушек Роговского. Гибридный подход, сочетающий точность измерений ТТ в нормальных режимах, с большей устойчивостью и отсутствием насыщения в переходных режима у КР, является перспективным направлением.

% DISS-005-B1380
Таким образом, использование катушек Роговского в качестве первичных преобразователей повышает надежность и эффективность релейной защиты и автоматики, обеспечивая стабильность электроснабжения промышленных объектов.\commentnote{DISS-005-C0094}{Алексей Евдаков [2]}{Хорошо бы добавить Мостик к ВНР. Вот только для него ток не очень важен, надо какого-то более общего характера этот мостик, без акцента на Катушках роговского.}

% DISS-005-B1381
Обеспечение точного и быстрого обнаружения аварийного режима и корректной работы измерительных органов, как это достигается с применением катушек Роговского, является первым и критически важным этапом функционирования систем автоматического ввода резерва. Однако, после успешного переключения на резервный источник питания, не менее значимой задачей становится безопасное и надежное восстановление нормальной схемы электроснабжения. Именно этому аспекту, а именно совершенствованию алгоритмов восстановления нормального режима (ВНР), посвящена следующая часть выводов, так как эффективность всего комплекса БАВР-ВНР определяется слаженной работой обеих его компонент.

% DISS-005-B1382

% DISS-005-B1383

% DISS-005-B1384

% DISS-005-B1385
\paragraph{На подумать}\commentnote{DISS-005-C0095}{Алексей Евдаков [2]}{Прост пометки, чтобы не забыть.}

% DISS-005-B1386
У разных производителей устройств БАВР, орган определяющий направление протекания тока (РНМ) различается, поэтому разбор существующих базовых проблем до введения адаптаций, и современных работ будет произведён на основе устройств производства ООО «АПС». Таким образом за «базовое исполнение РНМ» принимается логика, используемая в терминалах МИР производства ООО «АПС». А сравнительный анализ и оценка влияния описанных проблем будет проведена в конце раздела.

% DISS-005-B1387
Здесь речь идёт только про РНМ, но в моих анализах достаточно много и про пуск по углу. И работа в целом про БАВР, а не только РНМ. Однако идеи, связанные с внутренней логикой, могут считать корпоративной тайной (пока что). Вероятно, сиё может быть освещено только с точки зрения активации каких-либо функций и наличия самого факта проблем. Собственно, как и описывалось про РНМ.

% DISS-005-B1388
Пока что в данный разбор добавлять что либо, кроме РНМ не стал.

% DISS-005-B1389

% DISS-005-B1390
\paragraph{ИЗУЧИТЬ}

% DISS-005-B1391
В локальной библиотеке имеется папка «Библиотека\textbackslash Дисеры по адаптивным реле», требуется её изучить, очень вероятно, потом актуально будет делать ссылки на данный работы о том, что что-то по РНМ уже изучалось (если там РНМ)

% DISS-005-B1392
